\begin{table}[htbp]
\centering
\caption{Bands inside the 600--850\,\wn{} window that find no partner in the other library within 35\,\wn{}, expressed as a share of the intensity carried in that window.}
\label{tab:si-unmatched}
\small
\begin{tabular}{lrrrrr}
\toprule
Agent & Computed bands & Measured bands & Paired & Computed with no measured partner & Measured with no computed partner \\
\midrule
GD & 4 & 3 & 3 & 10\% & 0\% \\
GF & 4 & 4 & 3 & 50\% & 6\% \\
HN1 & 5 & 3 & 3 & 10\% & 0\% \\
HN2 & 2 & 0 & 0 & 100\% & 100\% \\
HN3 & 5 & 2 & 2 & 32\% & 0\% \\
A-230 & 7 & 6 & 6 & 8\% & 0\% \\
A-232 & 2 & 2 & 1 & 75\% & 88\% \\
A-234 & 2 & 2 & 2 & 0\% & 0\% \\
A-242 & 7 & 5 & 4 & 31\% & 38\% \\
\textbf{Mean} & \textbf{4.2} & \textbf{3.0} & \textbf{2.7} & \textbf{35\%} & \textbf{26\%} \\
\bottomrule
\end{tabular}

\begin{flushleft}\footnotesize Measured bands almost always find a computed partner: the share is zero for GD, HN1, HN3, A-230 and A-234. The asymmetry runs the other way. The calculation places on average 35\% of the window intensity in bands the measurement does not show, rising to 50\% for GF, which is one of the misassigned agents. The deficit in this window is therefore not that the calculation misplaces the measured bands but that it distributes intensity over a band set the measurement does not have. A-232 and HN2 are anomalous in both directions and are discussed separately in Sections~S5 and~S8.\end{flushleft}
\end{table}
